\documentclass[11pt,a4paper]{article}
\usepackage[T1,T2A]{fontenc}
\usepackage[utf8]{inputenc}
\usepackage[russian]{babel}
\usepackage[margin=2.1cm]{geometry}
\usepackage{amsmath,amssymb,amsthm}
\usepackage{enumitem}
\usepackage{booktabs,tabularx}
\usepackage{microtype}
\usepackage{xcolor}
\usepackage[colorlinks=true,linkcolor=blue!50!black,urlcolor=blue!50!black]{hyperref}
\setlist{itemsep=2pt,topsep=3pt}
\newcommand{\R}{\mathbb R}
\DeclareMathOperator{\Sym}{Sym}

\title{<<Решето>> для проблемы Арнольда 1971-4:\\ прототип программы, отсеивающей <<неподозрительные>> полиномы}
\author{}
\date{3 октября 2026 г.}

\begin{document}
\maketitle
\vspace{-6mm}

\section{Идея}
Почти каждый известный критерий неустойчивости --- это проверка \emph{полуалгебраического} условия на
конечный отрезок ряда Тейлора: <<форма $U_k$ принимает отрицательные значения>>, <<у $U_k$ нет критических
точек на сфере>>, <<эффективный потенциал на жёлобе имеет только простые нули>> и т.\,п. Поэтому для
конкретного полинома вопрос <<покрыт ли он критерием X?>> алгоритмически разрешим (Тарский--Зайденберг).
На практике его решают точной алгеброй (базисы Грёбнера, результанты, последовательности Штурма) или
численно. Программа применяет критерии по возрастанию стоимости и выдаёт вердикт с причиной. Всё, что
не покрыто ни одним реализованным критерием, она объявляет \emph{подозрительным}.

\section{Уровни решета (\texttt{sieve.py})}
\begin{center}\small
\begin{tabularx}{\textwidth}{>{\raggedright\arraybackslash}p{3.1cm}X}
\toprule
Уровень & Что проверяется \\
\midrule
0. Подготовка & $U(0)=0$, $\nabla U(0)=0$; эвристические тесты: изолирована ли критическая точка (не
является ли $\min_{S_r}|\nabla U|$ шумом округления при всех $r$) и есть ли отрицательные значения\\
1. Гессиан & отрицательное собственное значение $\to$ [Lyapunov]; ранг 3 $\to$ минимум; ранг 2 $\to$ [Pa20-T5]
(приведённый одномерный потенциал вычисляется рядом по лемме о расщеплении); ранг 1 $\to$ [Ko86] (форма на
ядре без минимума) или [Pa20-T5] (старшая форма приведённого потенциала без критических точек)\\
2. Первая форма $U_k$ & нечётная степень или отрицательные значения $\to$ [KP82]; плюс $\nabla U_k\neq0$ на
сфере $\to$ [Pa20-T3]; грани многогранника Ньютона с квазиоднородной частью без критических точек и малым
хвостом $\to$ [Pa20-T6]; $U_k>0$ $\to$ минимум\\
3. Жёсткое ядро: $U_k\ge0$ с нулями & разложение $U_k$ на множители; трассировка кривых нулей
(<<жёлобов>>) и поиск изолированных нулевых направлений; критерий [Bu25]; эффективные потенциалы на
жёлобах и вдоль нулевых направлений, с поправкой <<дополнения до квадрата>> за члены, линейные по
поперечной координате; проверка простоты нулей $\to$ [N13] (доказано в заметке о примере~13) или
[N13*] (тот же вириальный аргумент, доказательство только намечено)\\
4. Симметрии & 47 нетривиальных знаковых перестановок координат: инвариантная прямая или плоскость, на
которой $U$ не минимум, $\to$ сведение к $n=1$ или $n=2$ [SYM]\\
5. Вердикт & \texttt{TOTALLY\_UNSTABLE}, \texttt{UNSTABLE}, \texttt{MIN}, \texttt{NONISOLATED},
\texttt{SUSPICIOUS} (с причиной), \texttt{UNDECIDED} (вырожденный случай, но отрицательных значений не
найдено)\\
\bottomrule
\end{tabularx}
\end{center}
Для неустойчивости по Ляпунову достаточно одного <<хорошего рога>> --- компоненты множества $\{U<0\}$, на
которой работает вириальное поле: траектории с $H<0$ не могут перейти в другую компоненту. Для тотальной
неустойчивости нужны все рога.

\section{Проверка на примерах}
Тестовый набор (\texttt{tests.py}) включает примеры из основных заметок. Все вердикты ожидаемые:
$x^2+y^2-z^4$ $\to$ [Pa20-T5]; $x^2+y^4-z^4$ $\to$ [Ko86]; $x^4+y^4-z^6$ $\to$ [Pa20-T6] с весами
$(3,3,2)$; $(x^2+y^2-z^2)^2-z^6$ $\to$ [Bu25]; пример 13 $\to$ [N13]; пример 14 (жёлоб надо <<искать>>)
$\to$ \texttt{SUSPICIOUS}: эффективный потенциал тождественно зануляется после дополнения до квадрата;
$(x^2+y^2-z^2)^2+y^2z^4+\frac1{10}xyz^5-z^8$ $\to$ \texttt{SUSPICIOUS}: кратный нуль эффективного потенциала,
симметрий нет; $x^2y^2+y^2z^2+z^2x^2+xyz(x+y+z)^2$ $\to$ \texttt{NONISOLATED} (оси координат критические).

\section{Прогон на сгенерированных семействах}
\begin{center}\footnotesize\setlength{\tabcolsep}{4pt}
\begin{tabular}{llrrrrrrr}
\toprule
Сем. & Описание & $N$ & тот.\,н. & неуст. & мин. & неизол. & не реш. & подозр.\\
\midrule
A & общие разреженные, степени 2--6 & 200 & 10 & 186 & 0 & 1 & 0 & 3\\
B & $U_2=U_3=0$, степени 4--6 & 200 & 115 & 71 & 0 & 7 & 0 & 7\\
C & $(x^2+y^2-z^2)^2+W$, $\deg W=5..8$ & 400 & 117 & 70 & 1 & 1 & 4 & 207\\
D & квартика с кон. числом нулей $+W$ & 240 & 105 & 30 & 2 & 90 & 4 & 9\\
\bottomrule
\end{tabular}
\end{center}
Причины, по которым выживают <<подозрительные>>:
\begin{itemize}[leftmargin=*]
\item \textbf{A}: коранг 2, вырожденный приведённый потенциал --- 3.
\item \textbf{B}: жёлоб кратности $\ge4$ (не Морс--Ботт) --- 3; касающийся (чётнократный) нуль эфф. потенциала --- 2; нулевое направление с вырожденной поперечной частью --- 1; нечётнократный нуль эфф. потенциала --- 1.
\item \textbf{C}: касающийся (чётнократный) нуль эфф. потенциала --- 116; нечётнократный нуль эфф. потенциала --- 91.
\item \textbf{D}: нулевое направление с вырожденной поперечной частью --- 9.
\end{itemize}

\noindent\textbf{Комментарии.}
\begin{itemize}[leftmargin=*]
\item \emph{A и B (типичные полиномы):} 96--97\% отсеиваются уже на уровнях 1--2. Выжившие в A --- случаи
коранга 2 с приведённым потенциалом вида $-z^6(1+\dots)$: они почти наверняка покрываются сведением к
$n=2$ или взвешенным вириальным полем с нулевым весом по <<плоскому>> направлению.
\item \emph{C (класс жёлоба):} выживает 52\%, но это артефакт разреженной выборки. Моном
$x^ay^bz^c$ на окружности жёлоба даёт $\cos^a\varphi\,\sin^b\varphi$, то есть кратные нули при $a\ge2$ или
$b\ge2$. При плотных случайных коэффициентах эффективный потенциал имеет общий вид, и выживших почти нет.
Выжившие --- это <<рога>> с вырожденными кончиками (чётнократный нуль --- касание, нечётнократный ---
острый кончик). Их естественно атаковать следующим взвешенным раздутием в кончике.
\item \emph{D (конечное число нулей):} 90 случаев неизолированности настоящие: разреженное $W$ часто не
содержит мономов, которые <<шевелят>> луч нулей квартики, и тогда весь луч состоит из критических точек.
Выжившие имеют вид $x^4+y^4\pm z^5+{}$хвост. Теорема 6 Паламодова не срабатывает лишь потому, что
условие на хвост реализовано с запасом (взвешенная степень $\ge d+\max a_i$); при более точной форме
условия эти случаи, по-видимому, покрываются.
\item \emph{Ошибка, найденная по ходу.} Первый прогон показал, что моя проверка простоты нулей
пропускала нечётнократные нули вида $\cos^3\varphi$. Из-за этого 106 полиномов семейства C были ложно
<<доказаны>> критерием [N13]. После исправления их число упало с 217 до 111. Мораль: условия теорем нужно
проверять сертифицированно, иначе решето начинает <<доказывать>> лишнее.
\end{itemize}

\noindent\textbf{Динамическая проверка выживших.} Для трёх выживших запускались траектории с энергией $h<0$ из точек множества $\{U<0\}$ на сфере
$|q|=0.05$ (метод Рунге--Кутты 4-го порядка, шаг по локальной частоте). Траектория считается ушедшей,
если достигла сферы $|q|=0.3$.
\begin{center}\footnotesize\setlength{\tabcolsep}{4pt}
\begin{tabular}{p{9.2cm}ccc}
\toprule
Полином & траекторий & ушли & $\max|\Delta H|/|h|$\\
\midrule
S1: $(x^2+y^2-z^2)^2+y^2z^4+\frac1{10}xyz^5-z^8$ (касающийся нуль) & 36 & 36 & $2\cdot10^{-5}$\\
S2: $\bigl(x^2+y^2-z^2-z(x^2+xy)\bigr)^2-z^8$ (пример 14) & 60 & 60 & $7\cdot10^{-4}$\\
S3: выживший из C (тройной нуль эфф. потенциала)$^*$ & 24 & 24 & $2\cdot10^{-5}$\\
\bottomrule
\end{tabular}
\end{center}
{\footnotesize $^*$ $x^4+x^3y^2z^2-x^2y^3z+2x^2y^2-2x^2z^2+y^4-y^3z^3-2y^2z^2+z^4$; с грубым шагом дрейф энергии
был велик, поэтому приведён расчёт с шагом в 5 раз мельче.}

Как и ожидалось, выжившие ведут себя неустойчиво. Это подтверждает, что они <<подозрительны>> лишь
относительно реализованных критериев, а не как кандидаты в контрпримеры.


\section{Что важно понимать}
\begin{itemize}[leftmargin=*]
\item \emph{<<Подозрительный>> значит <<не покрыт реализованными критериями>>, а не <<кандидат в устойчивые>>.}
Каждый новый критерий съедает часть выживших. Критерий жёлоба из заметки о примере~13 появился именно так:
пример выжил, его разобрали, и это стало новым уровнем решета.
\item \emph{Случайный перебор почти бесполезен.} Типичные полиномы покрываются уже на уровнях 1--2.
Интересные случаи лежат на вырожденных стратах пространства коэффициентов (кратные нули эффективных
потенциалов, вырожденная поперечная структура и т.\,д.), поэтому их надо генерировать целенаправленно:
нормальные формы по $O(3)$ и масштабу плюс явно наложенные условия вырождения. Выжившие в разреженных
семействах выше --- как раз такие <<структурные>> вырождения.
\item \emph{Численные тесты эвристичны.} Для строгости нужны точная алгебра и интервальная арифметика.
Особенно это касается изолированности точки и поиска отрицательных значений в очень тонких рогах:
их ширина $\sim r^2$, $r^3$.
\item \emph{Устойчивость решето не докажет никогда.} Оно только сужает область, где может прятаться
контрпример. Перебор <<всех полиномов>> --- это полуалгоритм: бесконечный поток выживших, упорядоченный по
степени и высоте коэффициентов.
\item \emph{Метки [N13*]} опираются на тот же вириальный аргумент, что и доказанный случай [N13]
(жёлоб общего вида, рог вдоль нулевого направления, поправки дополнения до квадрата), но полностью я
доказал только случай кругового конуса.
\end{itemize}

\section{Как превратить прототип в инструмент}
\begin{enumerate}[leftmargin=*]
\item Итерированные взвешенные раздутия (шаг Ньютона--Пюизё вдоль жёлобов и в <<кончиках>> рогов).
Каждый вырожденный выживший порождает новую модельную задачу на следующем уровне.
\item Автоматический поиск вириальных полей $V$ с $\langle\nabla U,V\rangle\le\alpha U$ и $\Sym DV\ge-\alpha/2$
через линейное или полуопределённое программирование (SOS) в раздутых координатах.
\item Точная арифметика и сертифицированная численность: Singular или Macaulay2 для локальных алгебр и
факторизации, интервальные оценки для неравенств на сферах.
\item Перебор нормальных форм жёсткого ядра (квартики $Q^2$, $L^2P$, квартики с конечным числом нулей) с
наложением вырождений --- это и есть <<умный>> перебор подозрительных полиномов.
\end{enumerate}

\end{document}
